% !TeX root = ../main.tex
\section{Выводы}

В ходе выполнения работы проведено исследование и имитационное моделирование устройства обнаружения сигналов со стабилизацией уровня ложных тревог с оценкой меньшего среднего значения (ОБН с СУЛТ-МСЗ) для заданных параметров варианта ($q = 20\text{ дБ}$, $\text{ОПШ} = 15;\; 35\text{ дБ}$, $P_{\text{лт}} = 10^{-5}$, $N = 16$, $N_G = 1$, $N_M = 3$).

По результатам моделирования установлено, что в гомогенных условиях на фоне однородного собственного шума обнаружитель с СУЛТ-МСЗ обеспечивает стабилизацию вероятности ложной тревоги на заданном уровне $P_{\text{лт}} = 10^{-5}$ при значении порогового множителя $\alpha_{\text{МСЗ}} = 28.78$, а наличие защитных ячеек ($N_G = 1$) исключает эффект самомаскирования одиночной цели.

В гетерогенных условиях при наличии группы из $N_M = 3$ близкорасположенных мешающих целей, попадающих в одну половину скользящего окна ($N/2 = 8$ отсчетов), алгоритм МСЗ эффективно устраняет эффект взаимомаскирования. Оценка уровня шума формируется по свободной от мешающих сигналов половине окна, благодаря чему обеспечивается уверенное обнаружение близкорасположенных целей ($P_{\text{по}} = 0.7563$ при $\text{ОСШ} = 20\text{ дБ}$) с потерями $L = 3.96\text{ дБ}$ относительно оптимального обнаружителя Неймана--Пирсона.

При работе на границе ступенчатого перепада мощности помех ($\text{ОПШ} = 15\text{ дБ}$ и $35\text{ дБ}$) поведение обнаружителя носит выраженный асимметричный характер. До пересечения границы помех порог не завышается и сохраняет высокую вероятность обнаружения ($P_{\text{по}} \approx 0.76$) при стабильной вероятности ложной тревоги $P_{\text{лт}} < 10^{-5}$. Однако сразу после перехода тестируемого отсчета в область интенсивной помехи (пока отстающее полуокно остается в зоне слабого шума) выбор меньшего среднего приводит к сильному занижению порога и резкому всплеску вероятности ложной тревоги (до $P_{\text{лт}} \approx 0.41$ при $\text{ОПШ} = 15\text{ дБ}$ и $P_{\text{лт}} \approx 0.96$ при $\text{ОПШ} = 35\text{ дБ}$).

Таким образом, обнаружитель с СУЛТ-МСЗ целесообразно применять в радиолокационных системах для разрешения близкорасположенных целей и подавления точечных импульсных помех при условии их попадания не более чем в одну половину скользящего окна, тогда как для одновременной защиты от скачков протяженных помех требуется переход к комбинированным алгоритмам (например, СУЛТ-БЗ-ПС).

